% =====================================================================
% Macros : notations de l'énoncé et cadres « à venir »
% =====================================================================

% ---------- Notations (identiques à l'énoncé, au code et aux slides) ----------
\newcommand{\Fe}{F_e}
\newcommand{\Nfft}{N_{\mathrm{fft}}}
\newcommand{\Tw}{T_w}
\newcommand{\Tq}{T_q}
\newcommand{\SdB}{S_{\mathrm{dB}}}
\newcommand{\sthr}{s_{\mathrm{thr}}}
\newcommand{\gf}{\gamma_f}
\newcommand{\Nf}{N_f}
\newcommand{\Nt}{N_t}
\newcommand{\Nu}{N_u}
\newcommand{\nq}{n_q}
\newcommand{\Meff}{M_{\mathrm{eff}}}
\newcommand{\kw}{\kappa_w}
\newcommand{\mv}{m_v}
\newcommand{\kv}{k_v}
\newcommand{\dm}{\Delta m}
\newcommand{\dk}{\Delta k}
\newcommand{\RSB}{\ensuremath{\mathrm{RSB}}}
\newcommand{\TBD}{\ensuremath{\mathrm{TBD}}}
\newcommand{\TFD}{\ensuremath{\mathrm{TFD}}}
\newcommand{\hhat}{\hat{h}}
\newcommand{\E}{\mathbb{E}}
\newcommand{\Prob}{\mathbb{P}}

% ---------- Texte à compléter (rouge, bien visible) ----------
% Avant le rendu : chercher « À remplir » dans le PDF, il ne doit plus y en avoir.
\DeclareRobustCommand{\aremplir}[1]{\textcolor{red}{\textbf{[À remplir : #1]}}}

% ---------- Figures de A : cadre « figure à venir » tant que le PDF n'existe pas ----------
% Usage : \inclurefig{tbd_rsb}{ce que la figure montrera}
% Cherche <nom>.pdf puis <nom>.png, dans figures/ (Overleaf) puis ../figures/ (dépôt git).
% La figure est insérée à sa taille naturelle : MATLAB l'exporte déjà à sa largeur
% d'impression (utils/sauver_figure.m), donc son texte reste à 8-9 pt. Elle n'est
% réduite que si elle dépasse la largeur disponible. L'argument optionnel est ignoré
% (gardé pour compatibilité).
\newsavebox{\boitefig}
\newcommand{\figurenaturelle}[1]{%
  \sbox{\boitefig}{\includegraphics{#1}}%
  \ifdim\wd\boitefig>\linewidth\resizebox{\linewidth}{!}{\usebox{\boitefig}}\else\usebox{\boitefig}\fi}
\newcommand{\inclurefig}[3][]{%
  \IfFileExists{figures/#2.pdf}{\figurenaturelle{figures/#2.pdf}}{%
  \IfFileExists{figures/#2.png}{\figurenaturelle{figures/#2.png}}{%
  \IfFileExists{../figures/#2.pdf}{\figurenaturelle{../figures/#2.pdf}}{%
  \IfFileExists{../figures/#2.png}{\figurenaturelle{../figures/#2.png}}{%
  \fbox{\parbox[c][3.2cm][c]{0.92\linewidth}{\centering\small\color{gray}
    Figure à venir : #3\\[2pt] \texttt{\detokenize{#2.pdf}}}}}}}}}

% ---------- Tableaux de A (fichiers .tex booktabs, sans environnement table) ----------
% Usage : \inclureTableau{T02}{ce que le tableau contiendra}
\newcommand{\inclureTableau}[2]{%
  \IfFileExists{tableaux/#1.tex}{\input{tableaux/#1}}{%
  \IfFileExists{../resultats/tableaux/#1.tex}{\input{../resultats/tableaux/#1}}{%
  \fbox{\parbox[c][1.6cm][c]{0.92\linewidth}{\centering\small\color{gray}
    Tableau à venir : #2\\[2pt] \texttt{\detokenize{#1.tex}}}}}}}

% Case de tableau à remplir (version courte, pour ne pas élargir les tableaux)
\DeclareRobustCommand{\aq}{\textcolor{red}{\textbf{?}}}
